% from ega3-1-fr.tex
\oldpage[III]{109}
mêmes hypothèses de finitude sur $\mathcal S$ et $\mathcal M$) on conclut
de (\hyperref[III.2.4.1-fr]{2.4.1}) que
\[
\bigoplus_{k\in\mathbb Z}R^pf_*(\mathcal M_k)
\]
est un $\mathcal S$-Module de type fini pour tout $p$, et (sous les
hypothèses supplémentaires de (\hyperref[III.2.4.1-fr]{2.4.1}), (ii))
qu'il existe $N$ tel que pour $n\geq N$, on ait
$R^pf_*(\mathcal M_k(n))=0$ pour tout $p>0$ et tout $k\in\mathbb Z$, et
que l'homomorphisme canonique
$f^*(f_*(\mathcal M_k(n)))\to\mathcal M_k(n)$ soit surjectif pour tout
$k\in\mathbb Z$.
\end{env}

\subsection{Caractéristique d'Euler-Poincaré et polynôme de Hilbert}
\label{subsection:III.2.5-fr}

\begin{env}[2.5.1]
\label{III.2.5.1-fr}
Soient $A$ un anneau artinien, $X$ un $A$-schéma projectif sur
$Y=\operatorname{Spec}(A)$. Pour tout $\mathcal O_X$-Module cohérent
$\mathcal F$, les $H^i(X,\mathcal F)$ $(i\geq0)$ sont des $A$-modules de
type fini (\hyperref[III.2.2.1-fr]{2.2.1}), donc ici de longueur finie
puisque $A$ est artinien. On sait en outre
(\hyperref[III.2.2.1-fr]{2.2.1}) que $H^i(X,\mathcal F)=0$ sauf pour un
nombre fini de valeurs de $i\geq0$~; le nombre entier
\[
\chi_A(\mathcal F)=\sum_{i=0}^{\infty}(-1)^i
\operatorname{long}(H^i(X,\mathcal F))
\tag{2.5.1.1}\label{III.2.5.1.1-fr}
\]
est donc défini pour tout $\mathcal O_X$-Module cohérent $\mathcal F$.
Lorsque $A$ est un anneau local artinien, on dit que $\chi_A(\mathcal F)$
est la caractéristique d'Euler-Poincaré de $\mathcal F$ (par rapport à
l'anneau $A$). Pour $\mathcal F=\mathcal O_X$, on dit que
$\chi_A(\mathcal O_X)$ est le genre arithmétique de $X$ (par rapport à
$A$).
\end{env}

\begin{proposition}[2.5.2]
\label{III.2.5.2-fr}
Soit
\[
0\longrightarrow\mathcal F'\longrightarrow\mathcal F
\longrightarrow\mathcal F''\longrightarrow0
\]
une suite exacte de $\mathcal O_X$-Modules cohérents~; on a alors
\[
\chi_A(\mathcal F)=\chi_A(\mathcal F')+\chi_A(\mathcal F'').
\tag{2.5.2.1}\label{III.2.5.2.1-fr}
\]
\end{proposition}

\begin{proof}
Comme les modules de cohomologie de $\mathcal F$, $\mathcal F'$,
$\mathcal F''$ sont nuls sauf un nombre fini d'entre eux, il y a un
entier $r>0$ tel que la suite exacte de cohomologie s'écrive
\[
0\to H^0(X,\mathcal F')\to H^0(X,\mathcal F)
\to H^0(X,\mathcal F'')\to H^1(X,\mathcal F')\to\cdots
\]
\[
\cdots\to H^r(X,\mathcal F')\to H^r(X,\mathcal F)
\to H^r(X,\mathcal F'')\to0.
\]
Or, on sait que dans une suite exacte de $A$-modules de longueur finie,
ayant des 0 aux deux extrémités, la somme alternée des longueurs est
nulle (0, 11.10.1)~; appliquant ce résultat, on trouve immédiatement la
formule (\hyperref[III.2.5.2.1-fr]{2.5.2.1}).
\end{proof}

On notera que le résultat de (\hyperref[III.2.5.2-fr]{2.5.2}) s'applique
toutes les fois que l'on sait que $X$ est un $A$-schéma quasi-compact et
que les $A$-modules $H^i(X,\mathcal F)$ sont de type fini pour tout
$\mathcal O_X$-Module cohérent $\mathcal F$
(\hyperref[III.1.4.12-fr]{1.4.12}).

\begin{theorem}[2.5.3]
\label{III.2.5.3-fr}
Soient $A$ un anneau local artinien, $X$ un schéma projectif sur
$Y=\operatorname{Spec}(A)$, $\mathcal L$ un $\mathcal O_X$-Module
inversible très ample relativement à $Y$, $\mathcal F$ un
$\mathcal O_X$-Module cohérent~; on pose
$\mathcal F(n)=\mathcal F\otimes_{\mathcal O_X}\mathcal L^{\otimes n}$
pour tout $n\in\mathbb Z$.
\begin{enumerate}
\item[(i)] Il existe un polynôme unique $P\in\mathbb Q[T]$ tel que
$\chi_A(\mathcal F(n))=P(n)$ pour tout $n\in\mathbb Z$ (on dit que $P$
est le polynôme de Hilbert de $\mathcal F$ par rapport à $A$).
\item[(ii)] Pour $n$ assez grand, on a
$\chi_A(\mathcal F(n))=\operatorname{long}_A
\Gamma(X,\mathcal F(n))$.
\item[(iii)] Le terme de plus haut degré de $\chi_A(\mathcal F(n))$ a un
coefficient $\geq0$.
\end{enumerate}
\end{theorem}

Ajoutons qu'au chap. IV, dans le paragraphe consacré à la notion de
dimension, nous démontrerons en outre que le degré de
$\chi_A(\mathcal F(n))$ est égal à la dimension du support de
$\mathcal F$.

\begin{proof}
Comme on a $H^i(X,\mathcal F(n))=0$ pour tout $i>0$ dès que $n$ est
assez grand (\hyperref[III.2.2.1-fr]{2.2.1}),

